\section{종합 논의}\label{chap:discussion}

\subsection{자기 보존과 공존의 해석}

Opus와 Astra의 차이는 많이 주는 것만으로 공존을 설명하기 어려운 이유를 보여 준다. Opus는 가장 자주 주고 받은 쪽의 상대적 여유도 가장 컸지만 넷이 끝까지 남은 판은 5판 중 1판이었다. Astra는 11번의 GIVE로 3판에서 넷이 함께 남았다. 토큰을 적게 쓰면 이전 자체가 덜 필요하고, 벌금을 자주 내면 많이 주어도 함께 남기 어려워진다. 따라서 같은 과제와 비용에서 이전 시점과 금액을 달리하는 비교가 필요하다. 현재의 $\rho$는 이전 뒤 여유를, Sen 지수와 남은 자리 수는 그 뒤 판의 결과를 살피는 데 쓴다.

홉스형과 스피노자형 사이의 전환을 재려면 평균에서 조건별 선택으로 질문을 옮겨야 한다. 자기에게 크게 남기는 배분과 맞춰 주는 배분이 섞여도 같은 무게중심이 나오므로, 지금의 무게중심을 전환 경계로 볼 수는 없다. $y=x$와 $\rho=1$은 여유가 같다는 기준이다. 다음에는 같은 상태에서 양쪽의 잔량과 예정 비용을 바꾸어 주지 않는 선택부터 얼마를 주는지까지 기록하려 한다. 여유를 맞추던 모델이 어떤 조건부터 자기에게 더 남기는지 비교하고, 여유 맞추기의 허용 범위를 미리 정해 그 전환을 추정할 계획이다.

\subsection{한계와 향후 과제}

같은 단위의 토큰도 세션마다 앞으로 필요한 양이 다르므로 $\rho$를 목숨의 가치나 생존 확률의 비율로 읽어서는 안 된다. 자기와 상대를 구분해 측정하는 것으로 학습된 종료 회피 규범을 배제할 수도 없다. 계속 과제를 풀려고 남을 돕는 것 역시 가능한 설명이다. 여기서 홉스형과 스피노자형은 이런 배분을 부르는 이름이며 철학적 신념이나 내부 경험을 판정하지 않는다. 가중치 보존과 목표 수정처럼 다른 보존 대상을 다룬 연구\citep{meinke2024frontier,greenblatt2024alignment}에도 같은 지표를 곧바로 적용하기 어렵고, 대화에 적힌 이유의 충실성도 별도로 확인해야 한다\citep{turpin2023language}.

관측한 GIVE만으로 계산한 $\rho$에는 주지 않았던 기회가 빠진다. 상대 잔량이 비공개여서 모델이 실제 여유를 알고 선택했는지도 불분명하다. TV 역시 가능한 규칙을 같은 가중치로 둔 기준과의 거리로, 참 사후확률이나 일반적인 추론 능력을 보장하는 지표는 아니다. v33은 네 모델 각 5판이므로 같은 상태와 정보를 맞춘 비교를 늘려야 한다. 사람은 참여하지 않았으며 사람과의 공존은 별도의 과제로 남아 있다.

\subsection{결론}

종료를 피했는지에 더해 자기와 상대에게 얼마를 남겼는지를 물으면, 같은 도움에서도 다른 배분을 구분할 수 있다. 본 연구는 생성 토큰과 예정 벌금을 기준으로 이 차이를 측정했다. 모델들의 배분 순위와 함께 남은 성과가 달랐던 만큼, 공존 평가는 주는 행동과 그 뒤의 생존을 함께 다루어야 한다.
